% !TEX root = main.tex
% Codificación, idioma y tipografía
\usepackage[utf8]{inputenc}
\usepackage[T1]{fontenc}
\usepackage[spanish,es-tabla]{babel}
\usepackage{lmodern}
\usepackage{microtype}

% Márgenes solicitados para la plantilla.
% Ajustar únicamente si la universidad entrega medidas institucionales distintas.
\usepackage[
  letterpaper,
  top=2.5cm,
  bottom=2.5cm,
  left=3.0cm,
  right=2.5cm,
  headheight=15pt
]{geometry}

% Interlineado y párrafos
\usepackage{setspace}
\onehalfspacing
\setlength{\parindent}{1.25cm}
\setlength{\parskip}{0pt}

% Figuras, tablas y fórmulas
\usepackage{graphicx}
\graphicspath{{figuras/}}
\usepackage{booktabs}
\usepackage{longtable}
\usepackage{array}
\usepackage{tabularx}
\usepackage{ragged2e}
\usepackage{float}
\usepackage{pdflscape}
\usepackage{amsmath}
\usepackage{amssymb}
\usepackage{caption}
\usepackage{xcolor}
\usepackage{enumitem}

% Enlaces y referencias internas
\usepackage[hidelinks]{hyperref}
\hypersetup{
  pdftitle={Diseño de una arquitectura multiagente para la generación estructurada de artefactos de ingeniería de requisitos},
  pdfauthor={Carvajal Arias, Jhustyn Christopher},
  pdfsubject={Trabajo de titulación},
  pdfkeywords={ingeniería de requisitos, sistemas multiagente, RAG, trazabilidad}
}

% Ajustes generales
\setcounter{secnumdepth}{3}
\setcounter{tocdepth}{3}
\renewcommand{\bibname}{Referencias}
\renewcommand{\contentsname}{Índice general}
\renewcommand{\listfigurename}{Índice de figuras}
\renewcommand{\listtablename}{Índice de tablas}
\captionsetup{font=small,labelfont=bf}
\setlist[itemize]{topsep=3pt,itemsep=2pt,parsep=0pt}
\setlist[enumerate]{topsep=3pt,itemsep=2pt,parsep=0pt}
\newcolumntype{Y}{>{\RaggedRight\arraybackslash}X}
\newcommand{\codigo}[1]{\texttt{\detokenize{#1}}}
\newcommand{\pendiente}[1]{\textcolor{red!70!black}{[Pendiente: #1]}}
